%%
%% RGSW, External Product and Internal Product
%%
\section{TFHE Bootstrapping Primitives}
\begin{itemize}
    \item TFHE bootstrapping achieves fast evaluation of binary circuits. 
    \item It is based on (approximating) \textit{gadget decomposition}. 
    \item Asymmetry in noise scaling of GSW ciphertexts.
    \item Decomposition \cite{fhetextbook}: \textit{a mathematical technique used to represent a number in a smaller base (radix) while preserving its value} $\gamma = \gamma_0\frac{q}{\beta_0} + \dots + \gamma_\ell\frac{q}{\beta_\ell}$
    \item Gadget \cite{fhetextbook}: \textit{generalized form of number decomposition} $\gamma = \gamma_0g_0 + \dots + \gamma_\ell g_\ell$
\end{itemize}

\ifnotes
\begin{definition}[Generalized Gadget Matrix]
    Given a gadget vector $g$, the (generalized) gadget matrix $G_k$ \eqref{def:generalized_gadget_matrix} can be used to apply the gadget to each component of a vector of size $k$. \textcolor{red}{Generalizing here doesn't really serve a purpose cause always $k = 2$}
    
    \begin{equation}
        \label{def:generalized_gadget_matrix}
        G_k = I_k \otimes g = 
        \left(
        \begin{array}{c|c|c}
            g_0 & \dots & 0 \\
            \vdots & \ddots & \vdots \\
            g_{\ell - 1} & \dots & 0 \\
            \hline 
            \vdots & \ddots & \vdots \\ 
            \hline
            0 & \dots & g_0 \\
            \vdots & \ddots & \vdots \\
            0 & \dots & g_{\ell - 1} \\
        \end{array}
        \right)
        \in \mathcal{M}_{k\ell}
    \end{equation}    
\end{definition}
\fi

\begin{definition}[RGSW Ciphertext]
A (Ring) GSW ciphertext \eqref{def:rgsw}.. where $Z$ is a vector of $2\ell$ (fresh) \gls{rlwe} %$\rlwe_s(0)$ 
encryptions of 0. The message $m\in \{0,1\}$, and $G = I_2\otimes g$ where $g$ is the \textit{gadget vector}, typically defined as $g=(1, 1/B, \dots, 1/B^{\ell-1})$. 

\begin{equation}
    \label{def:rgsw}
    \rgsw_s(m) = Z + m\cdot G
\end{equation}
\end{definition}

\ifnotes
    \textcolor{red}{An \gls{rlwe} ciphertext is a vector $(b,a) = (a\cdot s + \Delta m + \epsilon)$ (and in scale-invariant schemes such as BGV, $\Delta = 1$). Typically the chosen gadget vector is $g = (1, B, \cdots, B^{\ell - 1})$. Therefore the RGSW ciphertext \eqref{def:rgsw} can be viewed as 2 sets of $\ell$ RLWE encryptions where the first set contains $\rlwe(mB^i)$ and the second set contains $\rlwe(-s\cdot mB^i)$ for $0 \leq i < \ell$.}
\fi

\begin{definition}[Gadget Property.]
    Two mappings $D(m)$ and $P(n)$ have the gadget property (wrt. $Q$) iff.
    \begin{equation}
        \label{property:gadget}
        \langle D(m), P(n)\rangle = m\cdot n \pmod{Q}
    \end{equation}
\end{definition}

The external product approximates the gadget property by defining a gadget matrix $G$ and inverse gadget matrix $G^{-1}(x)$ s.t.
\begin{equation}
    \label{def:approx_gadget_prop}
    G^{-1}(x)\cdot G = x + \epsilon \approx x \pmod{Q}
\end{equation}


\begin{definition}[External Product] $\boxdot$ defines a product between an \gls{rlwe} and an \gls{rgsw} ciphertext. $G^{-1}(x)$ satisfies the approximate gadget property \eqref{def:approx_gadget_prop}.
\begin{gather}
    \label{def:external_product}
    \mathsf{RLWE}_s(a) \boxdot \mathsf{RGSW}_s(b) \to \mathsf{RLWE}_s(a\cdot b)
    \\ a\boxdot B = G^{-1}(a)\cdot B \pmod{Q} 
\end{gather}
\end{definition}

\begin{definition}[Internal Product] $\boxtimes$ defines a product between \gls{rgsw} ciphertexts by repeated external products.
\begin{gather}
    \label{def:internal_product}
    \mathsf{RGSW}_s(a) \boxtimes \mathsf{RGSW}_s(b) = \mathsf{RGSW}_s(a\cdot b) \\
    % \nonumber \\ 
    A \boxtimes B = 
    \begin{bmatrix}
        A_0 \boxdot B \\
        \vdots \\
        A_{2\ell - 1} \boxdot B \\    
    \end{bmatrix} 
\end{gather}
\end{definition}

\subsection{Generalization}
The TFHE scheme \cite{cryptoeprint:2018/421} explicitly uses the gadget vector $g = (1, 1/B, \dots, 1/B^{\ell})$ \textcolor{red}{\texttt{Dec} algorithm}, but it should be noted that any choice of $G$ and $G^{-1}$ that satisfy the (inexact) gadget property \eqref{def:approx_gadget_prop} is sufficient to define a valid RGSW ciphertext. \ifnotes\textcolor{red}{TODO: establish this is true or cite, although it is trivial.}\fi 
For later it will be particularly helpful to establish constructively that/how any pair of vectors that satisfy the gadget property \eqref{property:gadget} can be used to find valid gadget matrices.

\ifnotes
\begin{gather}
    x\boxdot Y = G^{-1}(x)\cdot Y \\
    G^{-1}(x)\cdot (Z + mG) \\
    G^{-1}(x)Z + m(G^{-1}(x)\cdot G) \\
    (z + \epsilon_z) + m(x + \epsilon) \\
    \mathsf{RLWE}(0) + \mathsf{RLWE}(m\cdot m_x)
\end{gather}
\fi

\begin{theorem}%[Generalized Gadget RGSW]
    \label{the:gadget} % equal-length pair of vectors
    Any pair $D(a), P(b)$ that satisfies the approximate gadget property \eqref{def:approx_gadget_prop} defines an RGSW ciphertext \eqref{def:rgsw} and a valid external product \eqref{def:external_product} such that $\mathsf{RLWE}(a)\boxdot\mathsf{RGSW}(b) = \mathsf{RLWE}(a\cdot b)$. 
    Further, a valid (inverse) gadget matrix is given by:
    \begin{align}
        G &= I_2\otimes P(1)\\
        G^{-1}(x) &= [D(x_0)|D(x_1)]
    \end{align}
    
    \begin{proof}
        The inexact gadget property holds for $G= I_2\otimes P(1)$ and $G^{-1}(x) = [D(x_0)|D(x_1)]$.
        \begin{align*}
            G^{-1}(x)\cdot G &= [D(x_0)|D(x_1)] \begin{pmatrix}
                P(1) & 0 \\
                0 & P(1)
            \end{pmatrix} 
            \\ &= [D(x_0)\cdot P(1), D(x_1)\cdot P(1)] 
            \\ &= [x_0 \cdot 1 + \epsilon_0, x_1 \cdot 1 + \epsilon_1] 
            \\ &= [x_0, x_1] + [\epsilon_0, \epsilon_1] = x + \epsilon \pmod{Q}
        \end{align*}
        
        If the inexact gadget property holds, then the external product is valid.

        \begin{gather}
            G^{-1}(x)\cdot G = x + \epsilon_G \pmod{Q} \\
            x\boxdot Y = G^{-1}(x)\cdot (Z + m_y\cdot G) \\
            G^{-1}(x)\cdot Z + (G^{-1}(x)\cdot G)m_y \\
            G^{-1}(x)\cdot Z + (x + \epsilon_G)\cdot m_y \\
            G^{-1}(x)\cdot Z + x\cdot m_y + \epsilon_G \cdot m_y \\
        \end{gather}

        The middle term $x\cdot m_y = (a\cdot s + m_x + \epsilon, a)\cdot m_y = ((am_y)s + (m_ym_x) + m_y\epsilon, am_y)$. By setting $a' = a\cdot m_y, \epsilon' = \epsilon \cdot m_y$ its clear that this becomes a valid RLWE encryption of $m_y\cdot m_x$ from $(a's + m_x\cdot m_y + \epsilon', a')$.

        The left term $G^{-1}(x)\cdot Z$ evaluates to an RLWE encryption of 0 with some non-zero noise. (Here $g^{-1}_i$ denotes the $i$-th element of $G^{-1}(x)$, not necessarily the inverse of $g$...)
        \begin{gather}
            G^{-1}(x)\cdot Z = 
            \begin{pmatrix}
                g^{-1}_0 & \cdots &
                g^{-1}_{2\ell - 1}
            \end{pmatrix}
            \begin{pmatrix}
                b_0 & a_0 \\
                \vdots & \vdots \\
                b_{2\ell - 1} & a_{2\ell - 1} \\
            \end{pmatrix} \\
            =
            \left(
            \begin{array}{cc}
                \sum_{i=0}^{2\ell - 1}g^{-1}_ib_i, & \sum_{i=0}^{2\ell - 1}g^{-1}_ia_i 
            \end{array}
            \right)
        \end{gather}

        Let $a_z = \sum_{i=0}^{2\ell - 1}g^{-1}_ia_i$, and observe that the first element becomes
        \begin{gather}
            \sum_{i=0}^{2\ell - 1}g^{-1}_ib_i = \sum_{i=0}^{2\ell - 1}g^{-1}_i(a_i\cdot s + \epsilon_i) \\
            \sum_i g^{-1}_i a_i\cdot s + \sum_i g^{-1}_i e_i \\
            s\left(\sum_i g^{-1}_i a_i\right) + \sum_i g^{-1}_i e_i \\
            a_z\cdot s + \epsilon_z
        \end{gather}

        $\therefore \quad x\boxdot Y = \mathsf{RLWE}_s(0) + \mathsf{RLWE_s(m_x\cdot m_y) + (\epsilon_{G}\cdot m_y) = \mathsf{RLWE}_s(m_x\cdot m_y)}$
    \end{proof}
%\fi
\end{theorem}

%%
%% ADDED NOISE
\begin{remark}
    The exact gadget property is a special case of the inexact gadget property with $\epsilon = 0$, so the same proof holds. % Ideally we want to use a decomposition that has the \textit{exact} gadget property, to minimize noise. The noise added by the external product is also a function of the noise in $Z$, so this does not mean that the noise added by the external product is 0 even if $G, G^{-1}(x)$ has the exact gadget property. It also means some gadgets are \textit{better} than others, as they add less noise during the external product!
\end{remark}

\subsection{Noise}
Each term is a source of noise:
\begin{itemize}
    \item $\epsilon_z = \sum_i g^{-1}_i\epsilon_i$
    \item $\epsilon' = m_y\cdot \epsilon$
    \item $\epsilon_G\cdot m_y$
\end{itemize}

\noindent Some observations:
\begin{itemize}
    \item Two error terms are multiplied by $m_y$, which is why this method works best in binary systems where $m_y \in \{0, 1\}$ as this error term is bounded by the original error term $\|\epsilon'\| \leq \|\epsilon\|$
    \item The decomposition makes $\epsilon_z$ small because $g^{-1}_i$ are small and a sum of many small errors is better than a product of large errors (main point of decomposition). Also, since $\epsilon_i$ are sampled from a gaussian distribution, their expected value is $0$.
    \item The $\epsilon_G$ comes purely from the choice of gadget. If we choose an \textit{exact} gadget with $\epsilon = 0$ then $\epsilon_G\cdot m_y = 0$ and this error vanishes.
    \item $g^{-1}_i$ are also from the gadget, but is less controllable as it also depends on $x$. A good choice of gadget is therefore one that satisfies the exact gadget property and has small elements in $G^{-1}(x)$. 
\end{itemize}

The choice of gadget affects the noise and some gadgets are \textit{better} than others, as they add less noise during the external product! Ideally we want to use a decomposition that has the \textit{exact} gadget property, to minimize noise the noise from the \textit{``error (left) term``}.

% As we will see, we still need digit decomposition as in the original scheme to keep noise near constant, but the next sections will show how to do this in an RNS-native way (first decompose using the RNS primes, then decompose the towers).